\documentclass[11pt]{book}
\usepackage[paperwidth=145mm,paperheight=200mm,top=18mm,bottom=20mm,inner=20mm,outer=16mm]{geometry}
\usepackage{brumfiel}
% figures17.tex -- Appendix figures redrawn in TikZ.
%
% The Appendix (book pp.275-282) contains exactly ONE figure: the unnumbered
% and uncaptioned diagram on book p.277 that accompanies Postulate 16. It
% shows two three-ray configurations, unprimed and primed, drawn congruent.
% (Pages 275, 276, 278-281 are text only; 282 is blank. Verified page by page
% against sources/Geometry.pdf pp.289-296.)
%
% GEOMETRY MEASURED off sources/Geometry.pdf p.291 rasterised at 600 dpi.
% Method: threshold the ink, separate the heavy rays from the light arcs by
% binary erosion (rays survive a 5x5 erosion, arcs do not), take the vertex as
% the left end of the horizontal ray, then fit ray directions from the eroded
% ray cores only. Reading the whole ink mask instead biases every direction
% toward the arcs that run tangent to it -- that is what produced the earlier
% 45.0/117.0 figures. Raw and de-skewed (photo tilt removed by subtracting the
% measured tilt of the horizontal ray) results:
%
%     ray          raw unprimed  raw primed   de-skewed mean
%     horizontal       +0.48        +0.67          0 (datum)
%     diagonal         44.38        43.87        43.55
%     upper-left      116.33       116.13       115.66
%
% Drawn at 0 / 43.6 / 115.7 deg. Ray lengths and arc radii as fractions of the
% horizontal ray (mean of the two sub-figures):
%
%     horizontal ray 1.000   arc A 0.446   (drawn 3.30 / 1.47)
%     diagonal ray   0.994   arc C 0.699   (drawn 3.28 / 2.31)
%     upper-left ray 0.939   arc B 0.774   (drawn 3.10 / 2.55)
%
% The three arc radii are DISTINCT and ordered A < C < B, and each label sits
% on its own arc in a gap in that arc. Arc spans read off the source by
% scanning each arc's circle for ink (rays and label glyphs masked out):
%
%     arc A   4..60  and  78..114   gap 60..78, label at  69 deg
%     arc B  47..77  and  87..114   gap 77..87, label at  81.5 deg
%     arc C   4..16  and  30..42    gap 16..30, label at  23 deg
%
% The arrowhead ends below are the source's exactly; the GAP ends are opened a
% few degrees wider (A 52..86, B 71.5..91.5, C 12..34) because our labels are
% set with more padding than the book's tight hand lettering -- at the source's
% own gap widths check_labels.py reports the glyph boxes touching the arcs.
%
% Arrowheads therefore land at 114 (upper-left ray), 47/42 (diagonal ray) and
% 4 (horizontal ray) -- always just short of the ray they point at, and the two
% hits on the horizontal ray (arc A at 1.47, arc C at 2.31) and the two on the
% diagonal (arc B at 2.55, arc C at 2.31) are at different radii, as in the book.
%
% LINE WEIGHT: the book draws the three rays roughly 1.75x the weight of the
% annotation arcs (measured 13.0 px vs 8.5 px of ink at 600 dpi on the same
% photo). The rays keep the house `given' weight so the Appendix matches the
% rest of the book; the arcs are dropped to 0.42pt to restore the contrast.
%
% What the postulate rests on: the diagonal ray is INTERIOR to angle A, so
% angle A = angle B + angle C. See tools/constraints/ch17.py.

% ------------------------------------------------ Appendix figure (p.277)
\newcommand{\FIGAPPONE}{\begin{bkfigure}[0.52\linewidth]
% --- unprimed configuration
\begin{tikzpicture}[scale=0.72,
    mk/.style={fig, line width=0.42pt},          % light annotation arcs
    alab/.style={vlab, outer sep=1.5pt}]         % label sits in its arc's gap
  \coordinate (O) at (0,0);
  \draw[given] (O) -- (0:3.30);        % horizontal ray
  \draw[given] (O) -- (43.6:3.28);     % diagonal ray, interior to angle A
  \draw[given] (O) -- (115.7:3.10);    % upper-left ray
  % angle A : upper-left ray <-> horizontal ray  (the whole angle)
  \draw[mk,->] (86:1.47) arc[start angle=86, end angle=114, radius=1.47];
  \draw[mk,->] (52:1.47) arc[start angle=52, end angle=4,   radius=1.47];
  \node[alab] at (69:1.47) {$A$};
  % angle B : upper-left ray <-> diagonal ray
  \draw[mk,->] (91.5:2.55) arc[start angle=91.5, end angle=114, radius=2.55];
  \draw[mk,->] (71.5:2.55) arc[start angle=71.5, end angle=47,  radius=2.55];
  \node[alab] at (81.5:2.55) {$B$};
  % angle C : diagonal ray <-> horizontal ray
  \draw[mk,->] (34:2.31) arc[start angle=34, end angle=42, radius=2.31];
  \draw[mk,->] (12:2.31) arc[start angle=12, end angle=4,  radius=2.31];
  \node[alab] at (23:2.31) {$C$};
\end{tikzpicture}

\vspace{4mm}

% --- primed configuration (drawn congruent to the unprimed one, as the book
%     draws it; the postulate itself only supposes congruence of two pairs)
\begin{tikzpicture}[scale=0.72,
    mk/.style={fig, line width=0.42pt},
    alab/.style={vlab, outer sep=1.5pt}]
  \coordinate (O) at (0,0);
  \draw[given] (O) -- (0:3.30);
  \draw[given] (O) -- (43.6:3.28);
  \draw[given] (O) -- (115.7:3.10);
  \draw[mk,->] (86:1.47) arc[start angle=86, end angle=114, radius=1.47];
  \draw[mk,->] (52:1.47) arc[start angle=52, end angle=4,   radius=1.47];
  \node[alab] at (69:1.47) {$A'$};
  \draw[mk,->] (91.5:2.55) arc[start angle=91.5, end angle=114, radius=2.55];
  \draw[mk,->] (71.5:2.55) arc[start angle=71.5, end angle=47,  radius=2.55];
  \node[alab] at (81.5:2.55) {$B'$};
  \draw[mk,->] (34:2.31) arc[start angle=34, end angle=42, radius=2.31];
  \draw[mk,->] (12:2.31) arc[start angle=12, end angle=4,  radius=2.31];
  \node[alab] at (23:2.31) {$C'$};
\end{tikzpicture}
\end{bkfigure}}


% ---- local macros (hoist into book preamble) ----
% The Appendix has no numbered "N--M" sections. Its three major heads are
% set centred, bold and small, not flush left like a chapter's \section*.
% \APPhead also files the TOC entry the assembled book needs.
\newcommand{\APPhead}[2]{%
  \par\bigskip\begin{center}\bfseries #1\end{center}%
  \addcontentsline{toc}{section}{#2}\par\nopagebreak\medskip}
% Bold centred subhead grouping a run of postulates.
\newcommand{\APPsub}[1]{%
  \par\medskip\begin{center}\bfseries\small #1\end{center}\par\nopagebreak}
% Continuously-numbered reference list, with (a)/(b)/(c) subparts.
\newlist{applist}{enumerate}{2}
\setlist[applist,1]{label=\arabic*., leftmargin=1.7em, labelsep=0.4em,
                    itemsep=2.5pt, topsep=4pt, parsep=0pt}
\setlist[applist,2]{label=(\alph*), leftmargin=1.9em, labelsep=0.4em,
                    itemsep=1.5pt, topsep=2.5pt, parsep=0pt}
% Editorial summary standing in for a passage of the authors' connective
% prose that has NOT been reproduced verbatim -- see appendix-UNCERTAIN.md,
% item 1, which lists every one of these and the reason they are here.
\newcommand{\APPnote}[1]{%
  \par\medskip\noindent\fbox{\begin{minipage}{0.94\linewidth}\footnotesize
  \itshape #1\end{minipage}}\par\medskip}
% ---- end local macros ----

\begin{document}
\setcounter{chapter}{17}

% The book sets this opener flush RIGHT: "APPENDIX" occupies the slot that
% "CHAPTER 10" occupies on a numbered chapter opener (cf. book p.180).
\chapter*{\raggedleft\Huge APPENDIX}
\addcontentsline{toc}{chapter}{Appendix}
\markboth{APPENDIX}{}
\vspace{-1em}\noindent\rule{\linewidth}{0.6pt}\vspace{1em}
\figurekey

\APPnote{Opening paragraph (book p.~275): the Appendix collects all the
postulates for plane geometry together with a selection of the definitions
and theorems. The authors recommend reading it through rather than merely
consulting it, on the grounds that the order in which the theorems are
arranged is itself a picture of how Euclidean geometry is built up.
Summarised, not transcribed --- see \texttt{appendix-UNCERTAIN.md}, item~1.}

%========================================================= the postulates
\APPhead{THE POSTULATES}{The Postulates}

There are three groups of postulates, and three postulates listed singly.
They are:

\begin{center}
\begin{minipage}{0.86\linewidth}
Postulates of incidence (3 postulates)\\
Postulates of order or betweenness (5 postulates)\\
Postulates of congruence (7 postulates)\\
The postulate of Archimedes\\
The completeness postulate\\
The parallel postulate.
\end{minipage}
\end{center}

\APPsub{The Incidence Postulates}

\begin{applist}
\item There are at least three points not all on a line.
\item For any two different points, there is exactly one line containing
      these points.
\item Every line contains at least two points.
\end{applist}

\APPsub{The Betweenness Postulates}

\begin{applist}[resume]
\item If $B$ is between $A$ and $C$, then $A$, $B$, and $C$ are three
      different, or distinct, points on a line.
\item Of any three points on a line, exactly one is between the other two.
\item Four points on a line can be named $A_1, A_2, A_3, A_4$ so that the
      betweenness relations among them are exactly those given by the order
      of the subscripts.
\item If $A$ and $B$ are two points, then there is at least one point $C$
      such that $B$ is between $A$ and $C$, and at least one point $D$ such
      that $D$ is between $A$ and $B$.
\item Every line $l$ separates the plane: the points not on $l$ fall into
      two sets, the two \emph{sides} of $l$, such that
      \begin{applist}
      \item two points of the same set have no point of $l$ between them;
      \item two points of different sets have a point of $l$ between them.
      \end{applist}
\end{applist}

\APPsub{Linear Congruence Postulates}

\begin{applist}[resume]
\item Given $A$, $B$ on a line $l$ and $A'$ on a line $l'$, there is exactly
      one point $B'$ on a given side of $A'$ with $AB \cong A'B'$. Order is
      immaterial: $AB \cong A'B'$ also gives $AB \cong B'A'$,
      $BA \cong A'B'$ and $BA \cong B'A'$.
\item Segment congruence is reflexive, symmetric and transitive:
      \begin{applist}
      \item $AB \cong AB$.
      \item If $AB \cong A'B'$, then $A'B' \cong AB$.
      \item If $AB \cong A'B'$ and $A'B' \cong A''B''$, then
            $AB \cong A''B''$.
      \end{applist}
\item Segments may be added and subtracted. Let $B$ lie between $A$ and $C$
      on a line $l$, and $B'$ between $A'$ and $C'$ on a line $l'$.
      \begin{applist}
      \item If $AB \cong A'B'$ and $BC \cong B'C'$, then $AC \cong A'C'$.
      \item If $AC \cong A'C'$ and $BC \cong B'C'$, then $AB \cong A'B'$.
      \end{applist}
\end{applist}

\APPsub{Archimedes' Postulate}

\begin{applist}[resume]
\item Given a segment $AB$ on a line $l$ and any other segment $CD$, finitely
      many points $A_1, A_2, \ldots, A_n$ of $l$ can be laid off in order
      from $A$ with every segment $AA_1, A_1A_2, \ldots, A_{n-1}A_n$
      congruent to $CD$, and with either $B = A_n$ or $B$ between $A_{n-1}$
      and $A_n$.
\end{applist}

\APPsub{The Completeness Postulate}

\begin{applist}[resume]
\item For every line $l$, every point $A$ on $l$ and every positive real
      number $x$, there is a point $B$ on a given side of $A$ with
      $\sg{AB} = x$.
\end{applist}

\APPsub{Congruence Postulates for Angles}

\begin{applist}[resume]
\item If $\ang{A}$ is not a straight angle and $r'$ is a ray from $A'$ on a
      line $l'$, then on a given side of $l'$ there is exactly one ray $s'$
      from $A'$ making the angle at $A'$ with sides $r'$, $s'$ congruent to
      $\ang{A}$; we write $\ang{A'} \cong \ang{A}$. If $\ang{A}$ is a
      straight angle, the straight angle at $A'$ with side $r'$ is the only
      such angle.
\item Angle congruence is reflexive, symmetric and transitive:
      \begin{applist}
      \item $\ang{A} \cong \ang{A}$.
      \item If $\ang{A} \cong \ang{B}$, then $\ang{B} \cong \ang{A}$.
      \item If $\ang{A} \cong \ang{B}$ and $\ang{B} \cong \ang{C}$, then
            $\ang{A} \cong \ang{C}$.
      \end{applist}
\item In the figure, if any two of the three pairs $\ang{A}$ and
      $\ang{A'}$, $\ang{B}$ and $\ang{B'}$, $\ang{C}$ and $\ang{C'}$ are
      congruent, then so is the third pair. That is:
      \begin{applist}
      \item If $\ang{A} \cong \ang{A'}$ and $\ang{B} \cong \ang{B'}$, then
            $\ang{C} \cong \ang{C'}$.
      \item If $\ang{B} \cong \ang{B'}$ and $\ang{C} \cong \ang{C'}$, then
            $\ang{A} \cong \ang{A'}$.
      \item If $\ang{C} \cong \ang{C'}$ and $\ang{A} \cong \ang{A'}$, then
            $\ang{B} \cong \ang{B'}$.
      \end{applist}
      The angle $A$ may be a straight angle.
\end{applist}

\FIGAPPONE

\begin{applist}[resume]
\item If for $\tri{ABC}$ and $\tri{A'B'C'}$, $AB \cong A'B'$,
      $AC \cong A'C'$, and $\ang{A} \cong \ang{A'}$, then
      $\ang{B} \cong \ang{B'}$ and $\ang{C} \cong \ang{C'}$.
\end{applist}

\APPsub{The Parallel Postulate}

\begin{applist}[resume]
\item Through a point not on a line there is no more than one line parallel
      to the given line.
\end{applist}

%========================================================= definitions
\APPhead{SELECTED DEFINITIONS}{Selected Definitions}

\APPnote{Headnote (book p.~277): not every definition is equally
load-bearing, and the ones gathered here are the essential ones; a few are
given in abbreviated form where the short form cannot be misread.
Summarised, not transcribed.}

\begin{applist}
\item A \emph{geometric figure} is any set of points in the plane.
\item The \emph{line segment} (or \emph{segment}) determined by two points
      $A$, $B$ of a line is the set consisting of $A$, $B$ and all points
      between them, written ``$AB$''; $A$ and $B$ are its \emph{ends}.
\item For three points $A$, $B$, $C$ not all on a line, the points of the
      segments $AB$, $BC$, $AC$ form a \emph{triangle}, written
      $\tri{ABC}$; $A$, $B$, $C$ are its \emph{vertices} and $AB$, $AC$,
      $BC$ its \emph{sides}.
\item For $O$, $A$, $B$ on a line, $A$ and $B$ are on the \emph{same side of
      $O$} if $O$ is not between them, and on \emph{opposite sides of $O$}
      if $O$ is between them.
\item A \emph{ray} from $O$ consists of $O$ together with all points on one
      side of $O$ of a line through $O$.
\item An \emph{angle} is the set of all points on two rays from a point $O$,
      its \emph{vertex}; the rays are the \emph{sides}. If the rays lie on
      one line the angle is a \emph{straight angle}.
\item For an angle that is not straight, with vertex $O$ and sides $r_1$,
      $r_2$ on lines $l_1$, $l_2$: its \emph{interior} is the set of points
      lying both on the same side of $l_1$ as $r_2$ and on the same side of
      $l_2$ as $r_1$; its \emph{exterior} is the set of points belonging
      neither to the angle nor to its interior.
\item Two angles are \emph{adjacent} if they share a side and vertex and
      have no interior point in common.
\item Let $AB$ and $A'B'$ be segments, $A'B'$ on a line $l'$, and let $B''$
      be the point on the same side of $A'$ as $B'$ with $AB \cong A'B''$.
      Then $AB > A'B'$ if $B'$ is between $A'$ and $B''$, and
      $AB < A'B'$ if $B''$ is between $A'$ and $B'$.
\item For the definition of \emph{length of a segment} read Chapter 5,
      Section 4.
\item A \emph{right} angle is one congruent to its supplement.
\item Two triangles are \emph{congruent} if they can be labeled
      $\tri{ABC}$ and $\tri{A'B'C'}$ in such a way that all three pairs of
      corresponding sides and all three pairs of corresponding angles are
      congruent. This is written $\tri{ABC} \cong \tri{A'B'C'}$.
\item Let $\ang{A}$ and $\ang{B}$ be angles, neither of them straight. Copy
      $\ang{B}$ as $\ang{B'}$ against the first side of $\ang{A}$, placing
      it on the same side of that first side as the second side of
      $\ang{A}$. Then $\ang{B} < \ang{A}$ if the second side of $\ang{B'}$
      is interior to $\ang{A}$, and $\ang{B} > \ang{A}$ if it is exterior.
      If $\ang{A}$ is straight and $\ang{B}$ is not, then
      $\ang{A} > \ang{B}$.
\item Given $n$ points $P_1, P_2, \ldots, P_n$, the points lying on the
      segments $P_1P_2$, $P_2P_3$, \ldots, $P_{n-1}P_n$ form a
      \emph{polygonal path} joining $P_1$ to $P_n$; the points $P_i$ are
      its \emph{vertices}. If $P_1 = P_n$, the path is called a
      \emph{polygon}, designated polygon $P_1P_2 \ldots P_{n-1}$.

      A polygon is \emph{convex} if, for every two of its points, every
      point between them lies on the polygon or in its interior. A convex
      polygon that is both equiangular and equilateral is \emph{regular}.
\item Two lines are \emph{parallel} if they do not intersect.
\item If the ray $OB$ is in the interior of $\ang{AOC}$, or $\ang{AOC}$ is
      straight, then $\ang{AOC}$ is the \emph{sum} of $\ang{AOB}$ and
      $\ang{BOC}$.
\item The \emph{ratio} of segment $AB$ to segment $CD$ is the number
      $\sg{AB} \div \sg{CD}$. Segments $AB$, $CD$ are \emph{proportional} to
      $A'B'$, $C'D'$ when $\sg{AB}/\sg{A'B'} = \sg{CD}/\sg{C'D'}$; the sides
      of $\tri{ABC}$ are proportional to those of $\tri{A'B'C'}$ when
      $\sg{AB}/\sg{A'B'} = \sg{BC}/\sg{B'C'} = \sg{CA}/\sg{C'A'}$.
\item A \emph{circle} with center $O$ and positive number $r$ is the set of
      points $A$ with $\sg{OA} = r$. For $A$ on the circle, $OA$ is a
      \emph{radius} (plural \emph{radii}); a segment $AB$ with $A$, $B$ on
      the circle and $O$ on $AB$ is a \emph{diameter}.
\item For a circle with center $O$ and radius $r$: its \emph{interior} is
      the set of points $B$ with $\sg{OB} < r$, its \emph{exterior} the set
      of points $C$ with $\sg{OC} > r$. A \emph{chord} is a segment joining
      two points $E$, $F$ of the circle, and line $EF$ is a \emph{secant}. A
      \emph{central angle} has its vertex at the center; an \emph{inscribed
      angle} has its vertex on the circle and meets the circle again on each
      of its two sides.
\item If to each positive integer $n$ a number $x_n$ is assigned, the
      numbers $x_1, x_2, \ldots, x_n$ form an \emph{infinite sequence},
      written $\{x_n\}$, and $x_n$ is its $n$th term.
\item $L$ is the \emph{limit of the sequence} $x_n$ if $x_n$ is necessarily
      arbitrarily close to $L$ for large $n$; written $L = \lim \{x_n\}$.
\item The \emph{circumference} $C$ of a circle is the limit of the
      perimeters $p_n$ of the regular inscribed polygons: $C = \lim \{p_n\}$.
\item The number $C/d$, the same for every circle, is denoted by the Greek
      letter $\pi$ (pi).
\item The \emph{area} $S$ of a circle is the limit of the areas of the
      inscribed regular polygons: $S = \lim \{S_n\}$.
\item For a circular arc, let $l_n$ be the length of the regular polygonal
      path of $2^n$ sides inscribed in it. The \emph{length} $l$ of the arc
      is $l = \mathrm{limit}\, \{l_n\}$.
\item On a circle of radius $r$, an arc of length $\pi r/180$ is subtended
      by a central angle of \emph{one degree} ($1^\circ$).
\item If a central angle of a circle of radius $r$ subtends an arc of length
      $L$, then the \emph{measure of the angle in degrees} is
      $(L/\pi r) \cdot 180^\circ$.
\end{applist}

%========================================================= theorems
\APPhead{SELECTED THEOREMS}{Selected Theorems}

\APPnote{Headnote (book p.~280): the theorems below were picked to show how
the ideas were developed and, above all, \emph{in what order}; the list is
avowedly incomplete, and where a statement here is looser than the one in
the text the reader is asked to supply the missing precision or turn back to
the chapter. Summarised, not transcribed.}

\begin{applist}
\item Two different lines intersect in at most one point.
\item A point $O$ of a line $l$ separates $l$ into two \emph{sides} of $O$,
      such that for $P$, $Q$ on the same side either $OPQ$ or $OQP$, and for
      $P$, $Q$ on different sides $POQ$.
\item If $P$ and $Q$ are both interior to $\ang{O}$, so is the segment $PQ$.
\item A ray $r$ from $O$ other than the sides $r_1$, $r_2$ of $\ang{O}$ lies
      wholly in the interior or wholly in the exterior of $\ang{O}$.
\item If $P$ is interior to $\ang{O}$ and $Q$ exterior, the segment $PQ$
      contains a point of the angle.
\item For any two segments exactly one of $AB > A'B'$, $AB \cong A'B'$,
      $AB < A'B'$ holds.
\item If $AB > CD$ and $CD \cong EF$, then $AB > EF$.
\item If $AB \cong A'B'$, then $\sg{AB} = \sg{A'B'}$.
\item If $\sg{AB} = \sg{A'B'}$, then $AB \cong A'B'$.
\item $AB > CD$ if and only if $\sg{AB} > \sg{CD}$.
\item The S.A.S. theorem for congruent triangles.
\item The A.S.A. theorem.
\item The S.S.S. theorem.
\item For any two angles exactly one of $\ang{A} < \ang{B}$,
      $\ang{A} \cong \ang{B}$, $\ang{A} > \ang{B}$ is true.
\item All right angles are congruent to one another.
\item An exterior angle of a triangle is greater than either opposite
      interior angle.
\item In $\tri{ABC}$, $AB > AC$ if and only if $\ang{C} > \ang{B}$.
\item In $\tri{ABC}$, $\sg{AB} + \sg{BC} > \sg{AC}$.
\item If $l$ is any line and $A$ any point not on $l$, there exists at least
      one line on $A$ parallel to $l$.
\item If parallel lines are cut by a transversal, the alternate interior
      angles are congruent.
\item The sum of the angles of a triangle is a straight angle.
\item If a transversal is cut by three parallel lines in congruent segments,
      the same is true for every transversal of these lines.
\item If for $\tri{ABC}$ and $\tri{A'B'C'}$, $\ang{A} \cong \ang{A'}$,
      $\ang{B} \cong \ang{B'}$, and $\ang{C} \cong \ang{C'}$, then the
      triangles are similar.
\item If $\sg{AB}/\sg{A'B'} = \sg{AC}/\sg{A'C'} = \sg{BC}/\sg{B'C'}$, then
      $\tri{ABC}$ is similar to $\tri{A'B'C'}$.
\item In a right triangle the altitude on the hypotenuse forms two right
      triangles similar to the given triangle.
\item In $\tri{ABC}$ with $\ang{C}$ a right angle,
      $\sg{AC}^2 + \sg{BC}^2 = \sg{AB}^2$.
\item The area of a triangle is one half the base times the altitude.
\item Circles with centers $A$ and $B$, each of radius
      $r > \frac{1}{2}\sg{AB}$, intersect in two points on opposite sides of
      the line $AB$.
\item Exactly one circle is circumscribed about a regular polygon.
\item If regular polygons of $n$ sides are inscribed in two circles, then
      the perimeters of the polygons are proportional to the radii of the
      circles.
\item For a given circle let $p_n$ be the perimeter of a regular polygon of
      $n$ sides inscribed in the circle. The numbers $p_n$,
      $n = 3, 4, 5, \ldots$, form a sequence which has a limit.
\item The ratio $C/d$ (circumference to diameter) is the same for all
      circles.
\item If $S_n$ is the area of a regular polygon of $n$ sides inscribed in a
      circle, then $\mathrm{limit}\, \{S_n\}$ exists.
\item The area of a circle is equal to one half the circumference times the
      radius.
\item If two central angles are congruent, they subtend arcs of equal
      length.
\item If $\widehat{AB}$ and $\widehat{BC}$ are adjacent arcs (that is, if
      they have only point $B$ in common), then length $\widehat{ABC}$
      ${}={}$ length $\widehat{AB}$ ${}+{}$ length $\widehat{BC}$.
\item The measure of an angle inscribed in a circle is one half the measure
      of the central angle which subtends the same arc.
\end{applist}

\end{document}
